\chapter{Zeitlicher Verlauf der Sensoren und Aktoren}
\label{chap:zeitverlauf}

% Inhalt laut Vorgabe:
%  - Diagramm mit allen Sensoren und Aktoren über der Zeit
%  - 1. normaler Ablauf, mindestens 3 Fehlerfälle
%
% Hinweise zu tikz-timing:
%   H = high (true), L = low (false), D{Text} = Bus/Zustand mit Beschriftung
%   Zahl davor = Länge in Rastereinheiten, z. B. 3H = drei Einheiten high.
%   Jede Zeile muss in Summe gleich lang sein (hier 14 Einheiten = 7 s).

Die folgenden Diagramme zeigen alle Sensoren und Aktoren von Modul~1 sowie die
Handshake-Signale von Modul~2 über der Zeit. Ein Rasterkästchen entspricht
0,5~s. Hoch bedeutet \texttt{true}: bei Silo- und Trichtersensoren
\glqq erkannt\grqq, bei Aktoren \glqq an\grqq. Die beiden Not-Aus-Signale sind
sicherheitsgerichtet ausgelegt (Kapitel~\ref{sec:sensoren-aktoren}): Hoch
bedeutet \glqq Not-Aus nicht gedrückt\grqq, erst ein Wechsel auf low löst den
Stopp aus.

INIT und der direkt folgende Zustand BEREIT (\glqq BER.\grqq) dauern in der
Steuerung nur je einen Zyklus von 0,1~s. Sie sind zur besseren Lesbarkeit
mit einem Rasterkästchen dargestellt.

\section{Normaler Ablauf}
\label{sec:verlauf-normal}

Material liegt im Silo, Modul~2 ist frei. Die Steuerung startet die Druckluft,
schaltet nach 1~s das Förderband zu und fördert. Während der Förderung wird das
Silo kurzzeitig ganz gefüllt (Silo Top), anschließend leert es sich, bis auch
Silo Bottom kein Material mehr erkennt. Danach laufen Band und Luft 1~s nach.
Silo Top hat dabei keinen Einfluss auf den Ablauf.

Hinweis: Die Förderphase ist stark verkürzt dargestellt. In der realen Anlage
leert sich ein volles Silo nicht innerhalb von 2,5~s.

\begin{zeitdiagramm}{Zeitverlauf im normalen Ablauf}{fig:verlauf-normal}
  \resizebox{\textwidth}{!}{%
  \begin{tikztimingtable}[timing/slope=0, xscale=2.6, yscale=1.4, font=\small]
    Zustand        & 1D{INIT} 1D{BER.} 2D{VORLAUF} 5D{LÄUFT} 2D{NACHLAUF} 3D{BEREIT} \\
    Silo Bottom    & 9H 5L \\
    Silo Top       & 5L 2H 7L \\
    Trichter       & 14L \\
    Not-Aus        & 14H \\
    M2 Silo voll   & 14L \\
    M2 Not-Aus     & 14H \\
    Motor          & 4L 7H 3L \\
    Druckluft      & 2L 9H 3L \\
    Lampe rot      & 1H 13L \\
    Lampe orange   & 4H 5L 5H \\
    Lampe grün     & 1H 1L 9H 3L \\
  \extracode
    \tablerules
    \vertlines[help lines]{0,1,...,14}
  \end{tikztimingtable}}
\end{zeitdiagramm}

\clearpage
\section{Fehlerfall 1: Not-Aus während der Förderung}
\label{sec:verlauf-notaus}

Während LÄUFT wird der Not-Aus an Modul~1 gedrückt, das Signal wechselt auf
low. Motor und Druckluft gehen sofort aus, die rote Lampe leuchtet. Nach dem
Entriegeln ist das Signal wieder high, und der Ablauf startet erneut über
BEREIT und VORLAUF.

\begin{zeitdiagramm}{Zeitverlauf bei Not-Aus an Modul~1}{fig:verlauf-notaus}
  \resizebox{\textwidth}{!}{%
  \begin{tikztimingtable}[timing/slope=0, xscale=2.6, yscale=1.4, font=\small]
    Zustand        & 1D{INIT} 1D{BER.} 2D{VORLAUF} 4D{LÄUFT} 3D{STOP} 1D{BER.} 2D{VORLAUF} \\
    Silo Bottom    & 14H \\
    Silo Top       & 14L \\
    Trichter       & 14L \\
    Not-Aus        & 8H 3L 3H \\
    M2 Silo voll   & 14L \\
    M2 Not-Aus     & 14H \\
    Motor          & 4L 4H 6L \\
    Druckluft      & 2L 6H 4L 2H \\
    Lampe rot      & 1H 7L 3H 3L \\
    Lampe orange   & 4H 7L 3H \\
    Lampe grün     & 1H 1L 6H 4L 2H \\
  \extracode
    \tablerules
    \vertlines[help lines]{0,1,...,14}
  \end{tikztimingtable}}
\end{zeitdiagramm}

\clearpage
\section{Fehlerfall 2: Trichter voll}
\label{sec:verlauf-trichter}

Der Trichter am Schlauch meldet \glqq voll\grqq, zum Beispiel weil der Schlauch
verstopft ist. Die Steuerung geht in STOP, bis der Trichter wieder frei ist.

\begin{zeitdiagramm}{Zeitverlauf bei vollem Trichter}{fig:verlauf-trichter}
  \resizebox{\textwidth}{!}{%
  \begin{tikztimingtable}[timing/slope=0, xscale=2.6, yscale=1.4, font=\small]
    Zustand        & 1D{INIT} 1D{BER.} 2D{VORLAUF} 4D{LÄUFT} 3D{STOP} 1D{BER.} 2D{VORLAUF} \\
    Silo Bottom    & 14H \\
    Silo Top       & 14L \\
    Trichter       & 8L 3H 3L \\
    Not-Aus        & 14H \\
    M2 Silo voll   & 14L \\
    M2 Not-Aus     & 14H \\
    Motor          & 4L 4H 6L \\
    Druckluft      & 2L 6H 4L 2H \\
    Lampe rot      & 1H 7L 3H 3L \\
    Lampe orange   & 4H 7L 3H \\
    Lampe grün     & 1H 1L 6H 4L 2H \\
  \extracode
    \tablerules
    \vertlines[help lines]{0,1,...,14}
  \end{tikztimingtable}}
\end{zeitdiagramm}

\clearpage
\section{Fehlerfall 3: Modul 2 meldet Silo voll}
\label{sec:verlauf-m2-voll}

Während LÄUFT meldet Modul~2 über das Handshake-Topic \texttt{M2\_Silo\_Full},
dass sein Silo voll ist. Im Gegensatz zu den Fehlerfällen~1 und~2 ist das kein
Sicherheitsfehler. Modul~1 geht deshalb nicht in STOP, sondern beendet die
Förderung geregelt: Band und Druckluft laufen 1~s nach, damit der Schlauch leer
geblasen wird. Danach wartet Modul~1 in BEREIT (orange). Sobald Modul~2 wieder
Platz meldet, startet der Ablauf neu über VORLAUF.

\begin{zeitdiagramm}{Zeitverlauf bei vollem Silo an Modul~2}{fig:verlauf-m2-voll}
  \resizebox{\textwidth}{!}{%
  \begin{tikztimingtable}[timing/slope=0, xscale=2.6, yscale=1.4, font=\small]
    Zustand        & 1D{INIT} 1D{BER.} 2D{VORLAUF} 2D{LÄUFT} 2D{NACHLAUF} 2D{BEREIT} 2D{VORLAUF} 2D{LÄUFT} \\
    Silo Bottom    & 14H \\
    Silo Top       & 14L \\
    Trichter       & 14L \\
    Not-Aus        & 14H \\
    M2 Silo voll   & 6L 4H 4L \\
    M2 Not-Aus     & 14H \\
    Motor          & 4L 4H 4L 2H \\
    Druckluft      & 2L 6H 2L 4H \\
    Lampe rot      & 1H 13L \\
    Lampe orange   & 4H 2L 6H 2L \\
    Lampe grün     & 1H 1L 6H 2L 4H \\
  \extracode
    \tablerules
    \vertlines[help lines]{0,1,...,14}
  \end{tikztimingtable}}
\end{zeitdiagramm}
